% !TeX root = ../papers/nc-p3.tex
\section{階数15の倍率復元と一般添字保持同型類の次数}
\label{sec:rank-fifteen-parameter-orbits}
\begin{theorem}\label{thm:rank-fifteen-parameter-orbits}
正の固定曲線の第一表示の通常な階数15窓を取る。
節\ref{sec:rank-fifteen-pauli-comparison}の二つの積で生成する空間が十六次元なら、
倍率を任意に動かして同じ添字保持同型類を与える第一表示の全パラメータを、
以下の192個の射影線型変換と有理式で列挙できる。
幾何学的一般点では相異なる $(z,\lambda,n)$ はちょうど24個であり、
倍率 $(z,\lambda)$ はちょうど12組である。
第二表示を含めると同じ類を与える表示付きパラメータは48個になる。
\end{theorem}
\begin{proof}
$s^4=z$, $d^2=\lambda$, $t=sd$, $r=s^2$, $w=t^2$, $K=rk$, $N=rm/n$ とする。
二つの行列因子で、節\ref{sec:rank-fifteen-pauli-comparison}の平行移動生成元を
各平方が1になるよう規格化する。出力の $A,B$ を $r,\sqrt d$ で、入力のものを
$1,s\sqrt d$ でそれぞれ割る。各対は可換で正則な四文字平行移動として作用するから、
基底の一文字を選びその四つの像を取れば、対角基底変更で通常の平行移動になる。
符号行列はこの変更で変わらない。よってテンソルを
\[
 H=\sum_{x=(h,j)\ne0}\gamma_x P_x,\qquad
 P_x=X_h Z_j\otimes X_h Z_{\beta j}
\]
と書ける。ラベルの整数表示は $x=h+4j$、$\beta$ の像の列は $(0,3,2,1)$ である。
次の八量を定める。
\[
 \begin{gathered}
 a=1,\quad b=(K-1)/2,\quad L=n/2,\quad M=N/2,\\
 p=s+s^{-1},\quad q=s-s^{-1},\quad u=t+t^{-1},\quad v=t-t^{-1},\\
 A=(u-p)/4,\quad B=(u+p)/4,\quad C=(v-q)/4,\quad D=(v+q)/4 .
 \end{gathered}
\]
$\gamma_0=0$ を補うと全係数は次の順である。
\[
 (0,b,A,A,L,M,-C,D,\ a,a+b,B,B,L,M,-D,C).
\]
\path{computations/verify_rank_fifteen_parameter_orbits.py} は元のGram表示から
各単項行列へのトレース射影を求める。規格化前の行列 $P$ なら射影係数は
$\operatorname{tr}(HP)/(16P^2)$ である（$P^2$ はスカラー）。
$B$ を含む対ではさらに $sd$ を掛けて上記の規格化に戻す。全成分で差の分子が
四次式の分子で割り切れ、Gram逆行列は通常域で定義される。

共通スカラーも許したこれらの係数は
\[
 \begin{gathered}
 L^2+M^2=b(a+b),\qquad A^2+B^2=a(a+b),\qquad C^2+D^2=ab,\\
 LM=-AB=-CD
 \end{gathered}
\]
を満たす。逆に $a\ne0$ として $a=1$ に戻せば
\[
 s=B-A+D-C,\qquad t=B+A+D+C,\qquad n=2L
\]
と復元できる。実際 $s^{-1}=B-A-D+C$, $t^{-1}=B+A-D-C$ は上の二次式から従う。
$r=s^2,w=t^2$ とすれば $(r+1/r+w+1/w)/4=1+2b$,
$(r+1/r-w-1/w)/4=4LM$ である。従って $z=s^4$, $\lambda=(t/s)^2$ において
$K=1+2b$, $N=2M$ であり、第一式から $P(n)=0$ となる。
通常な階数15の候補だけを残せば、復元の両方向が成り立つ。

二つの局所交換子形式を保つ48個のラベル置換を $F$ とする。
$P_xP_y=(-1)^{e(x,y)}P_{x+y}$ の符号指数は $x=(h,j),y=(h',j')$ に対し
$e(x,y)=(j+\beta j)\cdot h'$ である。各 $F$ に対し、係数配列の変換は
$\gamma'_x=i^{v(x)}\gamma_{Fx}$（$i^2=-1$）と共通スカラーで記述できる。
指数は $P$ の平方と積で決まり、四基底ラベルで各二択が残る。二解の比はラベル群の
符号指標である。逆に各局所形式を保つ $F$ は規格化した四生成元の対応により
各行列因子の内部共役に持ち上がる。積の符号条件を満たす別の解は一方の因子の
Pauli行列の共役で得られる。従ってこれで全 $48\cdot16$ 変換を得る。
二積の生成次元が16なら、同型は全てこの形になる。

上記コードは八量を独立変数として、係数配列の線型関係を保つ変換を列挙する。
各 $F$ につき四つ、合計192個が残り、全てで五つの二次式がそれらの生成イデアルに入る。
これらは射影係数の表示を保つ群をなす。各 $F$ で他の変換も特殊点で有効になり得るかを確認する。
一つの保存変換を施した後の違いは符号指標だけである。
係数 $\gamma_2=\gamma_3$ と $\gamma_4=\gamma_{12}$ はともに非零なので、
指標は $1,8$ の両ラベルで値 $+1$ を取る必要がある。この条件を満たす指標は四つだけで、
既に列挙したものに尽きる。これで任意の生成次元16の点に対し、
192像から復元した通常点が必要十分なリストになる。生成次元を仮定しなくても、
このリストに入る点は添字保持同型である。

最後に $(z,\lambda,n)=(4,9,46)\bmod101$ で生成次元16は既知である。
$r=2,d=3$ とし、剰余二次拡大で $s^2=2$ を解く。コードで192変換を施すと、
全てが通常かつ非分岐で、相異なる $(s,t,n)$ が192個、$(z,\lambda,n)$ が24個、
倍率が12組となる。元の $n$ は単根である。等式と分母の非消滅検査により、
少なくとも24個の異なるパラメータが得られる条件は標数零の非空開集合で成り立つ。
逆に同型な各 $(z,\lambda,n)$ には $(s,d)$ の八通りの選択があり、
いずれも上の必要条件から192像に入る。従って上限は24である。
各倍率の第一表示で一般の一類を与える根は正負の一対であるため倍率数は12となる。
倍率反転が第一・第二表示を全単射で対応させるので、最後の個数も従う。
\end{proof}
